\section{Peta Jalur Ekspedisi: Taksonomi Materi Buku Ajar}

Demi merengkuh status gelar ksatria kelulusan sarjana utuh meramu pembedahan Algoritma *Computer Science*, ekspedisi mahasiswa di bangku kelas ini tidaklah cuma berhenti memandangi tabel kebenaran and/or tipis-tipis biasa \cite{rosen2019,levin_discrete}. 

Silabus buku ajar Mahakarya "Logika Informatika" ini dirancang berurut spiral layaknya arsitektur menara bertingkat berfondasi kokoh (Gambar~\ref{fig:peta-konsep}). Berikut adalah rute perbatasan negara (\textit{Roadmap}) yang bakal mahasiswa injang remuk pelan-pelan kelak:

\begin{figure}[htbp]
\centering
\begin{tikzpicture}[
  node distance=0.6cm and 1cm,
  box/.style={rectangle, draw=black, rounded corners, fill=blue!10, minimum width=2.4cm, minimum height=0.7cm, align=center, font=\small},
  root/.style={rectangle, draw=black, rounded corners, fill=green!20, minimum width=3.5cm, minimum height=0.8cm, align=center, font=\small\bfseries},
  arrow/.style={->, thick, >=stealth}
]
  \node[root] (root) {Kawah Candradimuka Logika Matematika};
  \node[box, below left=1.2cm and 0.5cm of root] (b3) {Bab 3 \& 4 \\Fondasi Proposisional Evaluatif};
  \node[box, below=1.2cm of root] (b4) {Bab 5 \& 6 \\ Aturan Inferensi Valid dan Sistem Pembuktian Matematis};
  \node[box, below right=1.2cm and 0.5cm of root] (b5) {Bab 7 \& 8\\Semesta Kuantifier Logika Predikat (SQL/AI)};
  \node[box, below=1.8cm of b4] (b6) {Bab 9 \& 10\\Teknik Mahkota Induksi \& Rekursi Canggih Algoritma};
  \draw[arrow] (root) -- (b3);
  \draw[arrow] (root) -- (b4);
  \draw[arrow] (root) -- (b5);
  \draw[arrow] (b3) -- (b4);
  \draw[arrow] (b4) -- (b5);
  \draw[arrow] (b5) -- (b6);
  \draw[arrow] (b4) -- (b6);
\end{tikzpicture}
\caption{Spiral Taksonomi Hirarki Logika - Masing-masing anak tangga mensuport pilar bab di masa atasnya.}
\label{fig:peta-konsep-bab2}
\end{figure}

\subsection{1. Pondasi Purba Biner: Bab III -- IV}
Ekspedisi start keras menjamah alam **Logika Proposisional** (Nilai Atomik Sempit) dan mematahkan anomali **Ekuivalensi/Tautologi** menggunakan artileri Tabel Kebenaran \textit{Brute-Force} serta pisau bedah 20 Hukum Aljabar Boolean De Morgan. (*Membentengi skill memangkas Bug percabangan sintaks If-Else Software*).

\subsection{2. Benteng Penarikan Konklusi Hakim: Bab V -- VI}
Mahasiswa dicecar latihan gila **Modus Ponens, Tollens,** dan \textbf{Pembangunan Proof Pembuktian Berundak}. Mendidik meracik kesimpulan 100\% tanpa celah memverifikasi struktur sistem tak patah didera skenario aneh. Metode Pembuktian Langsung (*Direct Proof*), Pembuktian Kontradiksi sesat, dan Contraposisi dihidangkan di piring hidangan panas.

\subsection{3. Naik Takhta Raja Variabel "Semua" \& "Sebagian": Bab VII -- VIII}
Bosan hanya mengupas si miskin entitas Atomik, Mahasiswa naik takhta mengeksplor luasnya dimensi alam semesta via tembak misil **Logika Predikat (Kuantor $\forall$/$\exists$)**. Ini fondasi nyawa perumusan Query (*WHERE EXISTS*) Tabel Relasional \textit{Database} dan pemodelan nalar berantai pada meskin Sistem Pakar \textit{Artificial Intelligence}.

\subsection{4. Seni Menyiasati Mesin Abadi: Bab IX -- X}
Pilar penutup agung membesut senjata rahasia \textbf{Logika Bukti Induksi Matematika} dan \textbf{Rentetan Rekursi (\textit{Recursion Loop})}. Mahasiswa diangkat statusnya mampu menjustifikasi membuktikan keabsahan \textit{Looping Algoritma Sorting/Searching} ber-triliun kali mutar *For-Loops* itu sah tak tersesat mati konyol membakar memori sirkuit!

Setiap rantai bab ini merupakan kunci sandi gembok masuk bagi gerbang bab urutan esok paginya. Tidak diperkenankan satupun sarjana meloncati gembok membolos \cite{levin_discrete,rosen2019}!
